\documentclass[12pt]{article}

%Laco e campo magnetico
%Faro, Fri Mar 14 19:22:09 WET 2014
%Written by Tordar

\usepackage{graphicx}
\usepackage{pst-all}
\usepackage{pst-3dplot}

\begin{document}

\pagestyle{empty}

\begin{pspicture}(-8,-8)(8,8)
%\psgrid
\psset{beginAngle=150,endAngle=60,linewidth=0.5mm,linecolor=black,linestyle=dashed}
%Laco:
\pstThreeDLine[linewidth=1.5mm,linecolor=gray,linestyle=solid,arrows=-](3,-2,0)(-3,2,0)(-3,2,3)(3,-2,3)(3,-2,0)
% Eixos:
\psset{linewidth=1mm,linecolor=black,linestyle=solid}
\pstThreeDLine[linewidth=1.5mm,arrows=->](0,0,0)(8,0,0)
\pstThreeDLine[linewidth=1.5mm,arrows=->](0,0,0)(0,0,8)
\pstThreeDLine[linewidth=1.5mm,arrows=->](0,0,0)(0,8,0)
\pstThreeDLine[linewidth=1mm,linestyle=dashed,dash=4mm 1mm,arrows=-](0,0,2)(0,3,2)
% Campo magnetico
\pstThreeDLine[linewidth=2mm,linecolor=blue,arrows=->](0,3,0)(0,6,0)
% Velocidade angular
\pstThreeDLine[linewidth=2mm,linecolor=blue,arrows=->](0,0,4)(0,0,7)
% Arco de rotacao da velocidade angular:
\pstThreeDCircle[beginAngle=-90,endAngle=180,linewidth=1mm,linecolor=red,arrows=->](0,0,5)(1,0,0)(0,1,0)
% Elemento de area:
\pstThreeDLine[linewidth=1.5mm,linecolor=blue,arrows=->](0,0,2)(2,3,2)
% Arco do angulo entre dS e B:
\pstThreeDCircle[beginAngle=57,endAngle=90,linewidth=1mm,linecolor=red,arrows=->](0,0,2)(2.5,0,0)(0,2.5,0)
% Texto:
\pstThreeDPut(8.5,0,0){\scalebox{1.5}{$X$}}
\pstThreeDPut(0,8.5,0){\scalebox{1.5}{$Y$}}
\pstThreeDPut(0,0,8.5){\scalebox{1.5}{$Z$}}
\pstThreeDPut(0,5,1){\scalebox{1.5}{$\mathbf{B}$}}
\pstThreeDPut(-3,2,3.5){\scalebox{1.5}{$S$}}
\pstThreeDPut(1,0,1.5){\scalebox{1.5}{$d\mathbf{S}$}}
\pstThreeDPut(1,0,7){\scalebox{1.5}{$\vomg$}}
\pstThreeDPut(-2,1,1){\scalebox{1.5}{$\phi$}}
\end{pspicture}

\end{document}
